O presente trabalho aborda a criação do RoadVet, uma aplicação móvel voltada para profissionais da área veterinária domiciliar registrarem prontuários, tutores e pets. A motivação para criar este aplicativo é explicada pelos problemas de inconsistência de rede presentes em áreas remotas, esses problemas comprometem a agilidade e segurança dos atendimentos clínicos dos animais de estimação, como também ocasionam a perda de dados importantes para os profissionais veterinários. Para desenvolver essa solução, foi utilizado o paradigma \textit{offline-first}, que fornece uma estrutura adequada para o sistema, permitindo que os dados sejam salvos na camada do cliente em casos de falha de conexão. Metodologicamente, o estudo caracteriza-se como uma pesquisa aplicada com desenvolvimento experimental. A arquitetura da solução se dividiu entre camada do cliente, implementada com a biblioteca \textit{React Native}, juntamente com o \textit{WatermelonDB} para a persistência local dos dados e a camada de suporte, que foi implementada com o uso do \textit{Django REST Framework} e o banco de dados relacional \textit{SQLite}. O processo de sincronização assíncrona foi projetado a partir do rastreamento das alterações locais, enviadas e recebidas em lote por meio de uma rota única da API, adotando a estratégia de resolução de conflitos \textit{Last-Write-Wins}. Os testes empíricos realizados em ambiente controlado comprovaram que o sistema é capaz de manter a operacionalidade completa do usuário sem conectividade e de sincronizar os dados de forma íntegra e automatizada após o restabelecimento da rede. Conclui-se que a abordagem offline-first mitiga os riscos de perda de dados e eleva a eficiência do atendimento veterinário volante.

\noindent\textbf{Palavras-chave}: Aplicação móvel. \textit{Offline-first}. \textit{React Native}. \textit{WatermelonDB}. Sincronização de Dados.

